\documentclass[11pt]{article}
\usepackage[margin=1in]{geometry}
\usepackage{amsmath, amssymb}
\usepackage{booktabs}
\usepackage{setspace}
\usepackage{enumitem}

\setstretch{1.2}
\setlength{\parindent}{0pt}
\setlist[itemize]{itemsep=0.4em}

\begin{document}

\begin{center}
{\Large \textbf{Solutions -- Tutorial: Inequality and Poverty in Canada}}\\
\vspace{1em}
{\large ECON 204 -- Contemporary Economic Issues}\\
\vspace{1em}
{\large Winter 2026}
\end{center}
\hrule
\vspace{1em}



%----------------------------------------------------
\section*{Question 1. Income Distribution and Inequality}

The annual before-tax household incomes (in Canadian dollars) are:
\[
(18{,}000,\ 22{,}000,\ 26{,}000,\ 30{,}000,\ 35{,}000,\ 42{,}000,\ 50{,}000,\ 65{,}000,\ 95{,}000,\ 180{,}000)
\]

\subsection*{(a) Mean and Median Income}

\textbf{Step 1: Total income.}
\[
18{,}000+22{,}000+26{,}000+30{,}000+35{,}000+42{,}000+50{,}000+65{,}000+95{,}000+180{,}000
=563{,}000
\]

\textbf{Step 2: Mean income.}
\[
\bar{y}=\frac{563{,}000}{10}=56{,}300
\]
The mean household income is \textbf{\$56,300}.

\textbf{Step 3: Median income.}  
With 10 households, the median is the average of the 5th and 6th incomes:
\[
\text{Median}=\frac{35{,}000+42{,}000}{2}=38{,}500
\]

\textbf{Interpretation.}  
The median income (\$38,500) is much lower than the mean. This occurs because very high incomes at the top of the distribution pull the mean upward. For this reason, the median is often a better measure of the ``typical'' household income in Canada.

\subsection*{(b) Income Shares}

\textbf{Step 1: Total income.}
\[
Y=563{,}000
\]

\textbf{Step 2: Bottom 40\% income share (households 1--4).}
\[
Y_{B40}=18{,}000+22{,}000+26{,}000+30{,}000=96{,}000
\]
\[
\text{Share}_{B40}=\frac{96{,}000}{563{,}000}\times 100 \approx 17.05\%
\]

\textbf{Step 3: Top 10\% income share (household 10).}
\[
Y_{T10}=180{,}000
\]
\[
\text{Share}_{T10}=\frac{180{,}000}{563{,}000}\times 100 \approx 31.97\%
\]

\textbf{Interpretation.}  
In this economy, the top 10\% of households earn nearly one-third of all income, while the bottom 40\% earn only about 17\%. This indicates substantial income concentration at the top.

%----------------------------------------------------
\section*{Question 2. Simple Inequality Ratios}

\subsection*{(a) Top-to-Bottom Ratio}

\[
\text{Top-to-bottom ratio}=\frac{180{,}000}{18{,}000}=10
\]

\textbf{Interpretation.}  
The richest household earns 10 times as much as the poorest household, illustrating a large income gap.

\subsection*{(b) Policy Change}

After the policy:
\[
\text{New ratio}=\frac{150{,}000}{25{,}000}=6
\]

\textbf{Conclusion.}  
The top-to-bottom ratio falls from 10 to 6, indicating that inequality decreases under the policy.

%----------------------------------------------------
\section*{Question 3. Poverty Measurement in Canada}

Poverty thresholds:
\begin{itemize}
\item LICO = \$26{,}000
\item MBM = \$28{,}000
\end{itemize}

\subsection*{(a) Poverty Status}

\textbf{Below LICO (strictly below \$26{,}000):}
\[
18{,}000,\ 22{,}000 \quad \Rightarrow \quad 2 \text{ households}
\]

\textbf{Below MBM (strictly below \$28{,}000):}
\[
18{,}000,\ 22{,}000,\ 26{,}000 \quad \Rightarrow \quad 3 \text{ households}
\]

\subsection*{(b) Poverty Rates}

\[
\text{LICO poverty rate}=\frac{2}{10}\times 100=20\%
\]
\[
\text{MBM poverty rate}=\frac{3}{10}\times 100=30\%
\]

\textbf{Interpretation.}  
The MBM poverty rate is higher because the MBM threshold is higher than the LICO threshold, capturing a more demanding standard of living.

%----------------------------------------------------
\section*{Question 4. Depth of Poverty}

MBM threshold = \$28{,}000.

\subsection*{(a) Poverty Gaps}

\[
\begin{aligned}
28{,}000-18{,}000 &= 10{,}000 \\
28{,}000-22{,}000 &= 6{,}000 \\
28{,}000-26{,}000 &= 2{,}000
\end{aligned}
\]

\subsection*{(b) Average Poverty Gap}

\[
\text{Average gap}=\frac{10{,}000+6{,}000+2{,}000}{3}=6{,}000
\]

\textbf{Interpretation.}  
The poverty gap measures how far households are below the poverty line, providing information about the severity of poverty, not just its incidence.

%----------------------------------------------------
\section*{Question 5. Inflation and Poverty}

\subsection*{(a) New MBM Threshold}

\[
28{,}000 \times 1.08 = 30{,}240
\]

The new MBM threshold is \textbf{\$30,240}.

\subsection*{(b) Poverty Impact}

Households below \$30,240:
\[
18{,}000,\ 22{,}000,\ 26{,}000,\ 30{,}000 \quad \Rightarrow \quad 4 \text{ households}
\]

\[
\text{New poverty rate}=\frac{4}{10}\times 100=40\%
\]

\textbf{Interpretation.}  
Inflation increases the cost of living and can raise poverty rates even if nominal incomes do not change.



\end{document}
